\chapter{Engineering a Real-Time Estimator in C++20}
\label{ch:engineering}

\begin{objectives}
\begin{itemize}
  \item Understand why a real-time estimator must be deterministic, and how the project's memory tiers achieve it.
  \item Use fixed-max-size Eigen matrices, member workspaces and allocation guards correctly.
  \item Avoid the Eigen pitfalls that silently corrupt numbers: dangling expression templates and aliasing.
  \item Design a test strategy in which every test can fail, and a systematic procedure for debugging a diverging filter.
\end{itemize}
\end{objectives}

The mathematics of the previous chapters is necessary but not sufficient. An estimator that is mathematically correct can still
miss its deadline because of an allocation, produce garbage because of an Eigen expression template, or pass all its tests because
the tests cannot fail. This chapter collects the engineering practices that the project encodes in AGENTS.md.

% ------------------------------------------------------------------------------
\section{Determinism and the memory tiers}

A flight controller consumes the estimator's output at a fixed rate. A late estimate is a wrong estimate. The main source of
unpredictable latency in C++ is dynamic memory allocation: \code{new}, a \code{std::vector} that grows, a dynamic-size Eigen temporary. The
allocator may take locks, search free lists or ask the operating system for pages, with latency that is usually small and
occasionally large. The project therefore separates three tiers (AGENTS.md E4).

\begin{center}\small
\begin{tabularx}{\textwidth}{@{}l X X@{}}
\toprule
Tier & Scope & Rule \\
\midrule
Load & parsing, building event arrays, constructing the filter & allocate freely, before streaming starts \\
Core & \code{vio\_core}, per IMU sample and per update & zero heap allocations \\
Frontend & OpenCV tracking & may allocate, must reuse buffers; allocations are measured and logged \\
\bottomrule
\end{tabularx}
\end{center}

\subsection{Fixed-max-size matrices}

The MSCKF covariance changes size as clones come and go, but never exceeds $15 + 6N_{max}$. Eigen supports exactly this case: a
matrix whose \emph{runtime} size is dynamic but whose \emph{maximum} size is a compile-time constant. Its storage is an inline array of the
maximum size, so resizing within the bound never allocates.

\begin{lstlisting}
inline constexpr int kMaxClones = 10;                  // decided in M10a
inline constexpr int kMaxDim    = 15 + 6 * kMaxClones;

using CovMatrix = Eigen::Matrix<double, Eigen::Dynamic, Eigen::Dynamic,
                                Eigen::ColMajor, kMaxDim, kMaxDim>;

CovMatrix P;          // inline storage for 75 x 75 doubles, no heap
P.resize(15, 15);     // allowed: within the compile-time maximum
\end{lstlisting}

Large workspaces ($S$, $K$, stacked $H$, temporaries for $\Phi P$) are class members, created once, rather than function-local
variables. Function-local fixed-max matrices would not allocate either, but they would live on the stack, and several of them at
once can exceed it (\cref{ch:cloning}).

\subsection{Proving it: allocation guards}

``It shouldn't allocate'' is not evidence. The project overrides the global \code{operator new} in its tests to count allocations
(\code{tests/support/alloc\_guard.hpp}), and wraps every core hot-path test with a guard that asserts a count of zero. In debug builds,
\code{EIGEN\_RUNTIME\_NO\_MALLOC} together with \code{Eigen::internal::set\_is\_malloc\_allowed(false)} makes Eigen itself abort on any internal
allocation and reports where it happened. The two mechanisms are complementary: the guard catches allocations anywhere, Eigen's flag
pinpoints those inside Eigen.

% ------------------------------------------------------------------------------
\section{Eigen pitfalls}

\paragraph{Expression templates and \code{auto}.} Eigen operations return lightweight \emph{expression} objects that are evaluated
lazily when assigned to a matrix. Storing one with \code{auto} keeps references to its operands, which may be temporaries that no
longer exist when the expression is finally evaluated.

\begin{lstlisting}
auto bad = R * (a + b);             // expression refers to the temporary (a + b)
Eigen::Vector3d x = bad;            // undefined behavior: the temporary is gone
Eigen::Vector3d good = R * (a + b); // evaluated immediately into a concrete type
\end{lstlisting}

The rule is simple: never use \code{auto} with Eigen expressions (AGENTS.md E5). Name the concrete type, or call \code{.eval()} when a
temporary really is wanted. AddressSanitizer catches many of these bugs, but not all.

\paragraph{Aliasing.} Eigen assumes that the destination of a matrix product may alias its operands and evaluates products into a
temporary. With dynamic sizes that temporary allocates. When the destination is known not to alias, \code{.noalias()} writes directly
into it. When it does alias, the temporary is required: write into a member workspace explicitly.

\begin{lstlisting}
tmp_15xN_.noalias() = Phi * P.topRightCorner(15, nc);  // reads P, writes workspace
P.topRightCorner(15, nc) = tmp_15xN_;                   // copy back
P.bottomLeftCorner(nc, 15) = tmp_15xN_.transpose();     // keep symmetry
\end{lstlisting}

Coefficient-wise operations do not alias-check. \code{A = A.transpose()} is the classic bug; use \code{A.transposeInPlace()} or symmetrize
with a workspace.

\paragraph{Passing and alignment.} Pass Eigen objects by \code{const\&} (or \code{Eigen::Ref} for blocks). Do not add
\code{EIGEN\_MAKE\_ALIGNED\_OPERATOR\_NEW}: since C++17, aligned \code{new} handles over-aligned types. The real alignment hazard is mixing
\code{-march} or alignment-changing flags between translation units or with dependencies, which the project forbids.

\paragraph{Decompositions without allocation.} Eigen's decompositions normally own their storage. Fixed-max decomposition objects held as
members, or the in-place variants that factor a matrix in its own storage, avoid per-call allocation. Whatever the API, verify with the
allocation guard.

% ------------------------------------------------------------------------------
\section{Time and events}

\begin{lstlisting}
using Clock = std::chrono::nanoseconds;
struct ImuSample   { Clock t; Eigen::Vector3d gyro, accel; };
struct ImageEvent  { Clock t; int index; };
using SensorEvent = std::variant<ImuSample, ImageEvent /*, PoseMeasurement */>;

// Delta t: subtract in integers, convert only the difference.
const double dt = std::chrono::duration<double>(t1 - t0).count();
\end{lstlisting}

Absolute time stays in 64-bit integer nanoseconds (\cref{sec:floating}). Events form one presorted sequence built at load time, so the
streaming loop is a single pass over contiguous memory: no priority queue, no allocation, no reordering surprises.

% ------------------------------------------------------------------------------
\section{Numerical robustness checklist}

\begin{itemize}
  \item Normalize every quaternion after every multiplicative update (propagation, injection, cloning).
  \item Symmetrize $P$ after propagation and after every update.
  \item Use the Joseph form; factor $S$ with Cholesky or LDL$^\mathsf{T}$; treat a failed factorization as an error to report, not to hide.
  \item Guard every $0/0$ with Taylor branches; use $\atantwo$ for angles; handle $\Log$ near $\pi$.
  \item Reject features with non-positive or tiny depth, poor parallax or large reprojection error before linearizing.
  \item Count and log gate rejections and triangulation failures.
  \item Never clamp negative eigenvalues of $P$ to hide a problem; find out why they appeared.
  \item In debug builds, assert finiteness at module boundaries so that a \code{NaN} is caught where it is born.
\end{itemize}

% ------------------------------------------------------------------------------
\section{Testing}

\begin{keyidea}[a test must be able to fail]
For every test, ask: would it still pass if I flipped the sign of X, transposed Y, or swapped two blocks? If yes, the test is
decoration. Run it once against a deliberately broken implementation to make sure it fails.
\end{keyidea}

The project's test strategy layers several kinds of evidence.
\begin{enumerate}
  \item \textbf{Algebraic unit tests}: identities such as $\Log(\Exp(\phi)) = \phi$ at zero, tiny, generic and near-$\pi$ angles.
  \item \textbf{Finite-difference Jacobian tests} built on $\bplus/\bminus$, at non-trivial states, with tolerances justified in a comment.
  \item \textbf{Closed-form trajectory tests}: synthetic trajectories with analytic IMU data and convergence-order checks.
  \item \textbf{Statistical tests}: Monte Carlo covariance checks and NEES/NIS with $\chi^2$ bounds, seeded with
        \code{std::mt19937\_64}. Because \code{std::normal\_distribution} differs between standard libraries, never hard-code values from one
        run; test statistical properties.
  \item \textbf{Invariance tests}: cloning enabled versus disabled must give identical IMU estimates when no visual updates occur.
  \item \textbf{Resource tests}: allocation guards around every core hot path.
  \item {\sloppy\textbf{Sanitizers}: AddressSanitizer and UndefinedBehaviorSanitizer, with
        \code{-fno-sanitize-recover=all} so that any report fails the test run.\par}
\end{enumerate}

% ------------------------------------------------------------------------------
\section{Debugging a diverging filter}

When the full system misbehaves, resist the urge to tune. Work down the layers instead.
\begin{enumerate}
  \item Run on \textbf{synthetic data} with exactly modeled noise. If it fails there, the bug is in the code.
  \item Disable visual updates. Does IMU-only propagation from ground truth behave as \cref{fig:cov-growth} predicts, with errors
        inside $3\sigma$?
  \item Replace visual updates by \textbf{pose updates} (Milestone 6). If those are consistent, the back end is sound.
  \item Feed the MSCKF \textbf{ground-truth feature tracks} from a synthetic scene. If that works, the problem is in the front end or the
        calibration.
  \item On real data, plot each residual component over time and look for bias, oscillation or outliers.
  \item Only then adjust noise parameters, and judge the result on consistency statistics.
\end{enumerate}

\section{Going further}

Once the baseline MSCKF runs, natural extensions are First-Estimates Jacobians (\cref{ch:observability}); online calibration of the
camera-IMU extrinsics and time offset~\cite{li2013msckf2, li2014timeoffset}; stereo observations from \code{cam1}; keeping long-lived
features as SLAM landmarks; zero-velocity updates; and evaluation on the harder EuRoC sequences (V1\_03, MH\_05) and other
datasets. Each extension adds states and Jacobians, and each Jacobian gets its derivation and its finite-difference test.

\begin{takeaways}
\begin{itemize}
  \item Determinism requires zero allocation in the core: fixed-max-size matrices, member workspaces, proven by allocation guards.
  \item Never use \code{auto} with Eigen expressions; use \code{.noalias()} only when there is no aliasing.
  \item Keep time in integer nanoseconds and process one presorted event stream.
  \item Every test must be able to fail; layer algebraic, Jacobian, closed-form, statistical and resource tests.
  \item Debug by removing layers until the failure disappears, not by tuning.
\end{itemize}
\end{takeaways}

\begin{exercises}
  \item Write a minimal allocation guard by overriding the global \code{operator new} and \code{operator delete}, and show that it detects a
        \code{std::vector::push\_back} beyond capacity but not a fixed-max Eigen \code{resize}.
  \item Construct an example in which \code{auto} with an Eigen expression produces wrong numbers without crashing. Does
        AddressSanitizer report it?
  \item Measure the time of 10\,000 covariance propagation steps with fixed-max and with fully dynamic Eigen matrices, in both build
        presets. Explain the difference.
  \item Take one of your Jacobian tests, break the Jacobian in three different ways, and record which breakages the test detects.
\end{exercises}
